% F11：本书原创矢量示意；依据所在节的定义、证明或明确模型。
\begin{figure}[htbp]
\centering
\begin{tikzpicture}[x=1cm,y=1cm,>=stealth,every node/.style={font=\small},line width=.65pt]
\draw[->,gray] (-0.12,0)--(5.4,0) node[right] {\(x\)};\draw[->,gray] (0,-0.12)--(0,2.8);
\node[] at (2.4,3.4) {窄峰函数};
\draw[AnalysisBlue,thick] plot coordinates {(4.2188,0.0000) (4.5000,2.5000) (4.7812,0.0000)};
\draw[AnalysisBlue,thick] plot coordinates {(2.1797,0.0000) (2.2500,1.2500) (2.3203,0.0000)};
\draw[AnalysisBlue,thick] plot coordinates {(1.1074,0.0000) (1.1250,0.6250) (1.1426,0.0000)};
\draw[AnalysisBlue,thick] plot coordinates {(0.5581,0.0000) (0.5625,0.3125) (0.5669,0.0000)};
\draw[AnalysisBlue,thick] plot coordinates {(0.2802,0.0000) (0.2812,0.1562) (0.2823,0.0000)};
\draw[AnalysisBlue,thick] plot coordinates {(0.1404,0.0000) (0.1406,0.0781) (0.1409,0.0000)};
\draw[->,gray] (6.88,0)--(12.1,0) node[right] {\(x\)};\draw[->,gray] (7,-0.12)--(7,2.8);
\node[] at (9.3,3.4) {峰点差商不趋于零};
\fill[orange!80!black](11.5,2)circle(1.7pt);\fill[orange!80!black](9.25,2)circle(1.7pt);\fill[orange!80!black](8.125,2)circle(1.7pt);\fill[orange!80!black](7.5625,2)circle(1.7pt);\fill[orange!80!black](7.28125,2)circle(1.7pt);\fill[orange!80!black](7.140625,2)circle(1.7pt);\draw[gray,dashed] plot coordinates {(7.0000,2.0000) (11.8000,2.0000)};
\node[] at (9.3,-0.6) {\(f(a_n)/a_n=1\)，坏集密度趋零};

\end{tikzpicture}
\caption[稀疏尖峰与近似微分]{稀疏尖峰与近似微分。使用本章习题的参数 \(a_n=2^{-n}\)、\(b_n=a_n^2/8\)。左图画出前六个三角峰；全部峰的支撑在原点密度为零，但峰点差商恒为一。}
\label{b1:fig:visual-f11}
\end{figure}
